% Einheit 5 von 5 – Prozentuale Veränderung (bank/prozentrechnung/e5.jsonl, zusammenbau v0.3)
%% TODO Zweigzeile Teil 1: was hier gelernt wird (Satz in Schülersprache aus dem Einheitstitel) – Einheit 5
\einheitenkopf[e5]{Einheit 5 von 5 · Prozentuale Veränderung}
\zweigzeile{ab Kl. 7, je nach Buch bis Kl. 8 · P10 oft}
\setcounter{aufgabe}{22}

%% TODO Titel: Ich-kann-Satz fehlt in der Bank (Platzhalter: Kettenname) – Nr. 23
%% TODO Anweisung: ein Satz über der ersten Teilaufgabe fehlt in der Bank – Nr. 23
\begin{aufgabe}{Welcher Wert ist der alte?}
\begin{teile}
\teil Der Preis stieg von $2{,}40\,€$ auf $2{,}70\,€$. – Welcher Wert ist der alte?
\end{teile}
\end{aufgabe}

%% TODO Titel: Ich-kann-Satz fehlt in der Bank (Platzhalter: Kettenname) – Nr. 24
%% TODO Anweisung: ein Satz über der ersten Teilaufgabe fehlt in der Bank – Nr. 24
\begin{aufgabe}{Veränderung}
%% TODO \rechenplatz{n}: Schrittzahl des Grundfalls fehlt in der Bank (nur bei fester Schreibform, 2.3 b)
\begin{teile}
\teil um $20\,\%$ gesenkt – was bedeutet dasselbe? Kreuze an.\\ \kreuz{auf $20\,\%$ gesenkt}\\ \kreuz{auf $80\,\%$ gesenkt}\\ \kreuz{um $80\,\%$ gesenkt}
\teil $80\,€$ um $10\,\%$ erhöht – neuer Preis? \leerfeld[€]
\teil $45\,€$ um $20\,\%$ erhöht – mit Faktor \leerfeld[€]
\teil von $40$ auf $50$ – um wie viel Prozent mehr? \leerfeld[\%]
\teil Nach $20\,\%$ Rabatt kostet ein Rad $360\,€$ – alter Preis? \leerfeld[€]
\teil netto (ohne Steuer) $200\,€$, $19\,\%$ Mehrwertsteuer – brutto? \leerfeld[€]
\teil Wahlbeteiligung $58\,\%$, später $64\,\%$ – wie viele Prozentpunkte mehr? Wie viel Prozent mehr? Runde auf eine Stelle.
\teil Eine Straße steigt auf $200$ m waagerechter Strecke um $14$ m – Steigung in Prozent? \leerfeld[\%]
\teil Ein Liter Milch kostete $1{,}15\,€$, jetzt kostet er $1{,}29\,€$. Um wie viel Prozent ist der Preis gestiegen? Runde auf eine Stelle. \hfill (P10 2026 FOR)
\teil Ein Busticket kostet $2{,}40\,€$ und wird um $25\,\%$ teurer. Wie viel kostet es dann? \hfill (P10 2024 OS) \\ \leerfeld[€]
\teil Kletterhalle: Erwachsene $14\,€$, Jugendliche (bis $17$ Jahre) $9\,€$. Familie Demir: Mutter, Vater, Tochter ($15$ Jahre), Sohn ($12$ Jahre). Montags gibt es $15\,\%$ Rabatt. Wie viel zahlt die Familie am Montag? \hfill (P10 2015 OS)
\teil Für Radwege gilt: Steigung höchstens $5\,\%$. Ergänze: „$5\,\%$ Steigung heißt: auf \leerfeld[m] waagerechter Strecke \leerfeld[m] Höhenunterschied.“ Herr Lenz sagt: „Ein Radweg, der auf $240$ m waagerechter Strecke um $15$ m ansteigt, ist zu steil.“ Hat er recht? Rechne. \hfill (P10 2025 OS)
\end{teile}
\end{aufgabe}

%% TODO Titel: Ich-kann-Satz fehlt in der Bank (Platzhalter: Kettenname) – Nr. 25
%% TODO Anweisung: ein Satz über der ersten Teilaufgabe fehlt in der Bank – Nr. 25
\begin{aufgabe}{Veränderung – Fehler finden, Begründen, Anwendung}
\begin{teile}
\teil Preis von $60\,€$ auf $75\,€$ – Steigerung in Prozent? Ida: \rechnung{15 : 75 &= 0{,}2 = 20\,\%} Fehler? Wie heißt es richtig?
\teil Warum bezieht man eine Preiserhöhung auf den alten Preis?
\teil Zwei Angebote für ein Fahrrad zu $600\,€$: A gibt $10\,\%$ Rabatt und auf den neuen Preis nochmals $10\,\%$; B gibt einmal $20\,\%$ Rabatt. Welches Angebot ist günstiger?
\end{teile}
\end{aufgabe}
